
\section{Introduzione}
Questo studio esamina se l'ospitare i Giochi Olimpici influisce significativamente 
sul numero di medaglie vinte dai paesi ospitanti. Abbiamo analizzato i 16 paesi 
che hanno ospitato almeno una volta tra il 1924 e il 2024, per un totale di 384 
partecipazioni olimpiche. 
\subsection{Ipotesi iniziale}
H0 (Ipotesi Nulla): Ospitare le olimpiadi non influisce sul numero di medaglie 
vinte da un paese.\\
H1 (Ipotesi Alternativa): Ospitare le olimpiadi influisce significativamente sul
numero di medaglie vinte.

\section{Metodologie}
\subsection{I dati}
I dati sono stati estratti da un dataset storico completo delle Olimpiadi moderne
presente su Kaggle, integrato dal 2017 in poi dai dati presenti su Wikipedia
(1924-2024), contenente informazioni su:
Paese partecipante (NOC - National Olympic Committee)
Anno dell'evento olimpico
Status di ospitalità binario
Numero totale di medaglie vinte (oro, argento, bronzo sommate)

\subsection{Il metodo}
Il test chi-quadrato è stato applicato a tutti i paesi con frequenze attese $\ge$ 
5. Per i paesi che non soddisfacevano questa assunzione (Grecia e Messico), i 
risultati sono riportati con un avviso di cautela (*), in quanto non 
statisticamente affidabili secondo gli standard di Cochran (1954).\\
Il test chi-quadrato di Pearson è stato utilizzato per testare l'indipendenza 
tra due variabili categoriche: ospitalità (sì/no) e performance medagliistica 
(alto/basso).

Formula:\\
$O_i$ = frequenza osservata\\
$E_i$ = frequenza attesa sotto ipotesi nulla
$$\chi^2= \sum_{i=1}^{2} \frac{(O_i - E_i)^2}{E_i}$$

Frequenze attese:
$$E_{ospita}=\text{Medaglie Totali} \times \frac{n_{ospita}}{n_{totale}}$$\\
$$E_{\text {non\_ospita}}=\text{Medaglie Totali} \times \frac{n_{non\_ospita}}
{n_{totale}}$$

\section{Risultati}
Quattordici paesi su sedici hanno mostrato un effetto ospitalità statisticamente 
significativo (p < 0.05). Tuttavia, due paesi (Grecia e Messico) non soddisfacevano 
le assunzioni del test chi-quadrato a causa di partecipazioni limitate e basso 
numero di medaglie totali (frequenze attese = 3.12 < 5). Questi risultati sono 
stati esclusi dall'analisi principale per non compromettere la validità statistica 
delle conclusioni.

\newcolumntype{C}{>{\centering\arraybackslash}X}

\begin{table}[H]
\centering
\caption{Analisi e risultati per ogni nazione}
\label{tab:medaglie_chi2}
\small 
\setlength{\tabcolsep}{2pt} 
\begin{tabularx}{\textwidth}{@{} l *{9}{C} @{}}
\toprule
\textbf{Paese} & \textbf{Nr. host} & \textbf{Nr. non host} & \textbf{tot medaglie} 
& \textbf{Med host} & \textbf{Med attese host} & \textbf{Med non host} & 
\textbf{Med attese non host} & \textbf{P Value} & \textbf{Risultati} \\
\midrule
AUS & 2 & 22 & 586  & 93  & 48.83  & 493  & 537.17 & 0.0000 & True  \\
BRA & 1 & 22 & 166  & 19  & 7.22   & 147  & 158.78 & 0.0000 & True  \\
CHN & 1 & 13 & 722  & 100 & 51.57  & 622  & 670.43 & 0.0000 & True  \\
ESP & 1 & 22 & 180  & 22  & 7.83   & 158  & 172.17 & 0.0000 & True  \\
GBR & 2 & 22 & 716  & 92  & 59.67  & 624  & 656.33 & 0.0000 & True  \\
FRA & 2 & 22 & 635  & 104 & 52.92  & 531  & 582.08 & 0.0000 & True  \\
GRE & 1 & 23 & 75   & 16  & 3.12*  & 59   & 71.88  & 0.0000 & True  \\
GER & 1 & 15 & 755  & 101 & 47.19  & 654  & 707.81 & 0.0000 & True  \\
USA & 3 & 20 & 2253 & 384 & 293.87 & 1869 & 1959.13& 0.0000 & True  \\
KOR & 1 & 18 & 315  & 33  & 16.58  & 282  & 298.42 & 0.0000 & True  \\
JPN & 2 & 20 & 538  & 87  & 48.91  & 451  & 489.09 & 0.0000 & True  \\
FIN & 1 & 22 & 243  & 22  & 10.57  & 221  & 232.43 & 0.0003 & True  \\
MEX & 1 & 23 & 75   & 9   & 3.12*  & 66   & 71.88  & 0.0007 & True  \\
NED & 1 & 23 & 341  & 23  & 14.21  & 318  & 326.79 & 0.0172 & True  \\
ITA & 1 & 23 & 630  & 36  & 26.25  & 594  & 603.75 & 0.0519 & False \\
CAN & 1 & 22 & 313  & 11  & 13.61  & 302  & 299.39 & 0.4696 & False \\
\bottomrule
\end{tabularx}
\vspace{2pt} 
\raggedright 
\scriptsize 
\textit{* = Questi risultati non soddisfano le assunzioni del test del $\chi^2$}
\end{table}

\subsection{Risultati significativi}
Risultati significativi con un p value di 0.00 sono AUS, BRA, CHN, ESP, GBR, FRA, 
GER, USA, KOR, JPN, FIN, per un totale di 11 stati. Per il NED, invece ha comunque
un valore molto basso di 0.0172.\\
Le frequenze attese di Grecia e Messico, per quanto diano un P value molto basso, 
come abbiamo detto non sono considerabili valide perchè non rispettano le assunzioni 
del test del $\chi$ quadro.

\subsection{Risultati non significativi}

ITALIA - P value: 0.0519\\
Per un totale di 24 partecipazioni, di cui in 1 ospita, per un totale di 630 medaglie \\
La frequenza attesa quando ospita: 26.25 e medaglie reali: 36, con una differenza 
di +9.75 medaglie. 
L'effetto positivo è debole e non statisticamente significativo, dimostrando un 
trend positivo (+36.9\%), tale che non possiamo escludere il caso con certezza statistica.\\
CANADA - P value: 0.4696\\
Per un totale di 23 partecipazioni, di cui in 1 ospita, per un totale di 313 medaglie\\
La frequenza attesa quando ospita: 13.61 e medaglie reali: 11, con una differenza 
di -2.61 medaglie.
Il trend negativo (-15.38\%) è l'unico riscontrato nei 16 stati analizzati, 
con il P value più alto trovato. 


\section{Conclusioni}
Dei 16 paesi analizzati, 2 (Grecia e Messico) sono stati esclusi per violazione delle
assunzioni del test, riducendo il campione valido a 14 paesi. Di questi, 12 
presentano una relazione statisticamente significativa tra ospitalità e performance sul 
medagliere (p < 0.05), suggerendo un effetto "home advantage" robusto nelle olimpiadi.
Italia e Canada sono gli unici due stati che non mostrano effetti significativi 
dall'ospitare i Giochi Olimpici.
Una limitazione dello studio riguarda i paesi con storia breve di ospitalità 
olimpica. Paesi come Grecia e Messico, che hanno ospitato solo una volta nel 
periodo analizzato, producono distribuzioni molto sbilanciate che violano le assunzioni 
parametriche. Studi futuri potrebbero utilizzare metodi non parametrici 
(es. Fisher's Exact Test) o raccogliere dati su periodi più lunghi per affrontare 
questo problema.